\label{blueprint:sec:lean}

A computer-assisted proof combines an analytic reduction with a finite
numerical calculation. Formalizing the periodic and open-chain spin-1
proofs of Refs.~\cite{OpenAI2026Periodic,OpenAI2026Boundaryfield}
directly in Lean~4~\cite{Moura2021Lean} faces three obstacles:
\begin{enumerate}
    \item \emph{Infinite-dimensional path space.} The continuous-time
          spatial transfer kernel of Ref.~\cite{OpenAI2026Periodic} acts on
          $L^2$ of Poisson bond histories (Sec.~\ref{obc:caps}); formalizing its
          trace identities directly would require path-space measure and operator
          machinery not yet in Mathlib.
    \item \emph{Model specificity.} Writing the analytic steps around the
          isotropic spin-1 Heisenberg Hamiltonian ties the formalization to one
          model, so none of the operator inequalities can be reused for other
          spin chains or ladders.
    \item \emph{Finite numerical inputs.} The seed bounds rely on external
          computations—block diagonalizations, short-ring trace recursions, and
          matrix-product trial contractions—whose floating-point or C/Python
          implementations cannot be trusted inside a proof assistant.
\end{enumerate}

The diagrammatic reformulation in the main text and
Secs.~\ref{sm:sec:trotter}--\ref{sm:sec:models} (and
Sec.~\ref{sm:sec:obc} for open chains) resolves all three
obstacles before formalization begins. First, discretizing imaginary
time into $M$ brickwork slices replaces the path-space kernel by the
finite-dimensional Gram matrix $\smT=\smF^{\mathsf T}\smF\succeq0$
[Eq.~\eqref{sm:eq:positive}]. Instead of constructing a limiting
continuous-time transfer operator in Lean, we prove the three closed
spatial power-sum inequalities~\eqref{sm:eq:closed} at finite $M$ and
pass them to the exact partition functions $Z_L(\beta)$ via the
fixed-$(L,\beta)$ matrix product limit of Sec.~\ref{sm:sec:limit}%
~\leanmark{https://sirui-lu.com/CertifiedGaps/docs/CertifiedGaps/Ring/Gap.html\#CertifiedGaps.RPRing.spatialMomentBounds}.
Likewise, to state ground-state uniqueness and
$\Delta_L=E_1(L)-E_0(L)>c$ without sorting the eigenvalues of $H_L$, we
require a single ground-state index $m_0$ separated from every other
level $m\ne m_0$ by $E_L(m)-E_L(m_0)>c$%
~\leanmark{https://sirui-lu.com/CertifiedGaps/docs/CertifiedGaps/PurityCriterion.html\#CertifiedGaps.GapGT.lt_sub_of_monotone}.
Second, isolating the bond reflection-positivity identity
$-h=\sum_\alpha C_\alpha\otimes FC_\alpha F$ with $FoF=o$
[Eq.~\eqref{models:eq:rp}] makes the periodic scale-doubling theorem
apply uniformly to any reflection-positive unit-cell ring%
~\leanmark{https://sirui-lu.com/CertifiedGaps/docs/CertifiedGaps/Ring/Gap.html\#CertifiedGaps.RPRing.gapGT_of_chainCertificate}.

\begin{table}[ht]
    \caption{Structure of the Lean~4 formalization~\cite{Lu2026CertifiedGaps}
        and its interactive proof blueprint~\cite{Lu2026Blueprint}. Each row
        lists a chapter of the blueprint, the corresponding manuscript
        sections, the theorem proved in Lean~4, and the finite numerical
        hypotheses (if any) on which that theorem depends.}
    \label{blueprint:tab:map}
    \footnotesize
    \begin{tabular}{@{}llll@{}}
        \toprule
        Blueprint chapter & Manuscript                                                           & Theorem proved in Lean~4 & Numerical hypotheses \\
        \midrule
        2.\ Two-purity criterion
                          & Main text; Sec.~\ref{sm:sec:gap}
                          & Uniform gap from purities $\ge1-\varepsilon_0$                                                                         %
        ~\leanmark{https://sirui-lu.com/CertifiedGaps/docs/CertifiedGaps/PurityCriterion.html\#CertifiedGaps.gapGT_of_seeds}
                          & None                                                                                                                   \\
        3.\ Asymmetric doubling \& chains
                          & Secs.~\ref{sm:sec:asymmetric}, \ref{models:sec:spectra}
                          & Uniform gap from rational doubling chain                                                                               %
        ~\leanmark{https://sirui-lu.com/CertifiedGaps/docs/CertifiedGaps/ChainCheck.html\#CertifiedGaps.gapGT_of_chainCertificate}
                          & None (kernel-checked)                                                                                                  \\
        4.\ Reflection-positive rings
                          & Secs.~\ref{sm:sec:transfer}--\ref{sm:sec:limit}, \ref{models:sec:rp}
                          & Reflection-positive ring gap theorem                                                                                   %
        ~\leanmark{https://sirui-lu.com/CertifiedGaps/docs/CertifiedGaps/Ring/Gap.html\#CertifiedGaps.RPRing.gapGT_of_chainCertificate}
                          & None                                                                                                                   \\
        5.\ Trace bounds \& sector tables
                          & Sec.~\ref{sm:sec:seeds}
                          & Seed purities from short-ring moments                                                                                  %
        ~\leanmark{https://sirui-lu.com/CertifiedGaps/docs/CertifiedGaps/Seed/Ring.html\#CertifiedGaps.RPRing.seeds_of_certificate}
                          & None                                                                                                                   \\
        6.\ Spin-1 Heisenberg ring
                          & Main text; Sec.~\ref{sm:sec:seeds}
                          & Spin-1 ring gap $\Delta_L>0.07290$ ($L\ge36$)                                                                          %
        ~\leanmark{https://sirui-lu.com/CertifiedGaps/docs/Leaderboard/Records/SpinOneHeisenberg.html\#Leaderboard.Records.SpinOneHeisenberg.temperature_lifting}
                          & Short-ring traces ($4\le n\le12$)                                                                                      \\
        7.\ Open chains with end fields
                          & Sec.~\ref{sm:sec:obc}
                          & Open spin-1 gap ($n\ge1921$)                                                                                           %
        ~\leanmark{https://sirui-lu.com/CertifiedGaps/docs/CertifiedGaps/Heisenberg/OpenChain.html\#CertifiedGaps.SpinOneHeisenberg.spin_one_boundary_field_gap}
                          & Four seeds at $\beta_0=49$                                                                                             \\
        8.\ Spin chains \& ladders
                          & Sec.~\ref{sm:sec:models}; Table~\ref{models:tab:summary}
                          & Uniform gaps for Table~\ref{tab:models} models                                                                         %
        ~\leanmark{https://sirui-lu.com/CertifiedGaps/docs/CertifiedGaps/Models/StaggeredHeisenberg.html\#CertifiedGaps.StaggeredHeisenberg.gap}
                          & Seed pair $(p_0,q_0)$                                                                                                  \\
        9.\ Kernel-checked seeds
                          & Secs.~\ref{sm:sec:seeds}, \ref{models:sec:ising}
                          & Unconditional Ising gap; $S=1$ trials                                                                                  %
        ~\leanmark{https://sirui-lu.com/CertifiedGaps/docs/SeedKernel/Ising/Seed.html\#SeedKernel.Ising.ising_gap}
                          & None (Ising); $n\ge5$ traces                                                                                           \\
        \bottomrule
    \end{tabular}
\end{table}

Third, separating the universal doubling recurrence from the finite seed
allows the formalization to treat the numerical inputs at two levels
(Table~\ref{blueprint:tab:map}):
\begin{enumerate}
    \item \emph{Seed-hypothesis theorems (Chapters~6--8).} Because the
          asymmetric update~\eqref{sm:eq:update} is a rational map and the
          interval condition~\eqref{models:eq:chain-cert} replaces logarithms by
          rational power comparisons $r^{2^j}z_j<1-z_j$, the step from any
          rational seed pair $(p_0,q_0)$ to a uniform gap $\Delta_L>(\log r)/\beta_0$
          reduces to finitely many rational inequalities checked in the Lean
          kernel. Taking $(p_0,q_0)$ (or the short-ring trace table of
          Table~\ref{cert:tab:traces}) as an explicit hypothesis instantiates
          machine-checked gap theorems for all eight periodic models in
          Table~\ref{tab:models} and for the open spin-1 chain with endpoint
          fields (Theorem~\ref{obc:theorem}), while allowing the seed enclosures
          to be certified externally by verified block diagonalizations
          (Sec.~\ref{models:sec:spectra}) or exact integer trace recursions
          (Sec.~\ref{sm:sec:seeds}).
    \item \emph{Kernel-checked seed arithmetic (Chapter~9).} Where the
          finite seed calculation can be reduced to exact integer identities, the
          Lean kernel verifies the seed calculation itself. For the mixed-field
          Ising chain at $(g,h)=(1/4,1/20)$ (Sec.~\ref{models:sec:ising}), the
          kernel checks the integer trial-vector residuals~\eqref{models:eq:residual}
          and the classical configuration count~\eqref{models:eq:ising-rho},
          proving the uniform gap~\eqref{models:eq:ising-gap} unconditionally
          without external numerical hypotheses%
          ~\leanmark{https://sirui-lu.com/CertifiedGaps/docs/SeedKernel/Ising/Seed.html\#SeedKernel.Ising.ising_gap}.
          For the periodic and open spin-1 chains, the kernel checks the exact
          integer contractions~\eqref{cert:eq:mps-contractions} of the $72$- and
          $120$-site periodic matrix-product trial states, the four open-chain
          trial states on $34$, $78$, $142$, and $240$ sites, and the shortest
          ring trace recurrences ($n=4$ periodic and nine of fifteen boundary
          traces), reducing the remaining hypotheses to the short-ring integer
          trace totals for $5\le n\le12$%
          ~\leanmark{https://sirui-lu.com/CertifiedGaps/docs/SeedKernel/Heisenberg/Integers.html\#CertifiedGaps.SpinOneHeisenberg.gap_of_hornerTotalsLarge}
          and the remaining six boundary traces%
          ~\leanmark{https://sirui-lu.com/CertifiedGaps/docs/SeedKernel/Heisenberg/Integers.html\#CertifiedGaps.SpinOneHeisenberg.spin_one_boundary_field_gap_of_hornerTotalsLarge}.
\end{enumerate}

The human-readable proof blueprint~\cite{Lu2026Blueprint}
(\url{https://sirui-lu.com/CertifiedGaps/blueprint/}) links every
definition, lemma, and theorem in the dependency graph to its Lean~4
declaration, and the companion repository~\cite{Lu2026CertifiedGaps}
bundles the manuscript, standalone numerical certificate checkers, and
formal proofs following the Agentic Publication Protocol~\cite{Lu2026APP}.
